\documentclass[10pt,a4paper]{amsart}
\usepackage{amsmath,amssymb,amsthm,graphicx}
% Retained preamble for the Stacks Project chapter morphisms.tex (Morphisms of Schemes).
% The chapter file pinned-original-source.tex carries only \input{preamble} on its first line, so no
% preamble text lives inside the chapter itself.  The Stacks Project uses one single shared preamble for
% every chapter; the snapshot stages it verbatim in sources/stacks-chow/preamble.tex.  The block below is
% copied mechanically, line for line, from that file's lines 14-18 (the Xy-pic setup) and lines 154-217
% (the native theorem declarations together with the mathematical macros).  Nothing was added, removed or
% redefined.  The wrapper class and the amsmath/amssymb/amsthm/graphicx base packages are supplied by the
% assembler, so the class/\externaldocument part of the shared preamble is not repeated here.
% For commutative diagrams we use Xy-pic
\usepackage[all]{xy}

% We use 2cell for 2-commutative diagrams.
\xyoption{2cell}
% Theorem environments.
%
\theoremstyle{plain}
\newtheorem{theorem}[subsection]{Theorem}
\newtheorem{proposition}[subsection]{Proposition}
\newtheorem{lemma}[subsection]{Lemma}

\theoremstyle{definition}
\newtheorem{definition}[subsection]{Definition}
\newtheorem{example}[subsection]{Example}
\newtheorem{exercise}[subsection]{Exercise}
\newtheorem{situation}[subsection]{Situation}

\theoremstyle{remark}
\newtheorem{remark}[subsection]{Remark}
\newtheorem{remarks}[subsection]{Remarks}

\numberwithin{equation}{subsection}

% Macros
%
\def\lim{\mathop{\mathrm{lim}}\nolimits}
\def\colim{\mathop{\mathrm{colim}}\nolimits}
\def\Spec{\mathop{\mathrm{Spec}}}
\def\Hom{\mathop{\mathrm{Hom}}\nolimits}
\def\Ext{\mathop{\mathrm{Ext}}\nolimits}
\def\SheafHom{\mathop{\mathcal{H}\!\mathit{om}}\nolimits}
\def\SheafExt{\mathop{\mathcal{E}\!\mathit{xt}}\nolimits}
\def\Sch{\mathit{Sch}}
\def\Mor{\mathop{\mathrm{Mor}}\nolimits}
\def\Ob{\mathop{\mathrm{Ob}}\nolimits}
\def\Sh{\mathop{\mathit{Sh}}\nolimits}
\def\NL{\mathop{N\!L}\nolimits}
\def\CH{\mathop{\mathrm{CH}}\nolimits}
\def\proetale{{pro\text{-}\acute{e}tale}}
\def\etale{{\acute{e}tale}}
\def\QCoh{\mathit{QCoh}}
\def\Ker{\mathop{\mathrm{Ker}}}
\def\Im{\mathop{\mathrm{Im}}}
\def\Coker{\mathop{\mathrm{Coker}}}
\def\Coim{\mathop{\mathrm{Coim}}}

% Boxtimes
%
\DeclareMathSymbol{\boxtimes}{\mathbin}{AMSa}{"02}

%
% Macros for moduli stacks/spaces
%
\def\QCohstack{\mathcal{QC}\!\mathit{oh}}
\def\Cohstack{\mathcal{C}\!\mathit{oh}}
\def\Spacesstack{\mathcal{S}\!\mathit{paces}}
\def\Quotfunctor{\mathrm{Quot}}
\def\Hilbfunctor{\mathrm{Hilb}}
\def\Curvesstack{\mathcal{C}\!\mathit{urves}}
\def\Polarizedstack{\mathcal{P}\!\mathit{olarized}}
\def\Complexesstack{\mathcal{C}\!\mathit{omplexes}}
% \Pic is the operator that assigns to X its picard group, usage \Pic(X)
% \Picardstack_{X/B} denotes the Picard stack of X over B
% \Picardfunctor_{X/B} denotes the Picard functor of X over B
\def\Pic{\mathop{\mathrm{Pic}}\nolimits}
\def\Picardstack{\mathcal{P}\!\mathit{ic}}
\def\Picardfunctor{\mathrm{Pic}}
\def\Deformationcategory{\mathcal{D}\!\mathit{ef}}

% --- appended 2026-10-12 (Chief): the project own definitions of the reference/slogan/history
% display environments, copied mechanically from the shared Stacks preamble lines 8-12
% (\usepackage{verbatim} plus the three \newenvironment lines, which the chapter preamble keeps
% verbatim from the project preamble). Nothing added, removed or redefined: these environments
% carry no mathematical content and the project itself typesets them as comments. Without these
% three lines a counted statement whose own text uses them fails with Environment slogan undefined.
\usepackage{verbatim}
\newenvironment{reference}{\comment}{\endcomment}
\newenvironment{slogan}{\comment}{\endcomment}
\newenvironment{history}{\comment}{\endcomment}
\usepackage{cleveref}
\crefname{originallemma}{original lemma}{original lemmas}
\Crefname{originallemma}{Original lemma}{Original lemmas}
\makeatletter
\newlabel{constructions-lemma-morphism-relative-proj}{{07ZG}{source}{}{sourcepointer-1}{}}
\newlabel{constructions-lemma-morphism-relative-proj@cref}{{[originallemma][0][]07ZG}{[1][0][]source}{}{}{}}
\makeatother
\crefname{originallemma}{original lemma}{original lemmas}
\Crefname{originallemma}{Original lemma}{Original lemmas}
\makeatletter
\newlabel{constructions-lemma-surjective-generated-degree-1-map-relative-proj}{{07ZK}{source}{}{sourcepointer-2}{}}
\newlabel{constructions-lemma-surjective-generated-degree-1-map-relative-proj@cref}{{[originallemma][0][]07ZK}{[1][0][]source}{}{}{}}
\makeatother
\crefname{originalsection}{original section}{original sections}
\Crefname{originalsection}{Original section}{Original sections}
\makeatletter
\newlabel{constructions-section-projective-bundle}{{01OA}{source}{}{sourcepointer-3}{}}
\newlabel{constructions-section-projective-bundle@cref}{{[originalsection][0][]01OA}{[1][0][]source}{}{}{}}
\makeatother
\crefname{originallemma}{original lemma}{original lemmas}
\Crefname{originallemma}{Original lemma}{Original lemmas}
\makeatletter
\newlabel{lemma-relatively-very-ample-separated}{{01VQ}{source}{}{sourcepointer-4}{}}
\newlabel{lemma-relatively-very-ample-separated@cref}{{[originallemma][0][]01VQ}{[1][0][]source}{}{}{}}
\makeatother
\crefname{originallemma}{original lemma}{original lemmas}
\Crefname{originallemma}{Original lemma}{Original lemmas}
\makeatletter
\newlabel{schemes-lemma-push-forward-quasi-coherent}{{01LC}{source}{}{sourcepointer-5}{}}
\newlabel{schemes-lemma-push-forward-quasi-coherent@cref}{{[originallemma][0][]01LC}{[1][0][]source}{}{}{}}
\makeatother
\title{Four equivalent descriptions of a relatively very ample invertible sheaf for a quasi-compact morphism}
\author{The Stacks Project Authors}
\date{stacks-project:morphisms.tex:local-snapshot-20261012; GNU FDL 1.2 or later}
\begin{document}
\maketitle
\noindent\textit{Editorial note.} All delivered mathematical text is the author's own, copied from the pinned Stacks Project chapter file pinned-original-source.tex; no statement or proof step was rewritten, reordered or removed.\par\medskip
\noindent\textit{Editorial note.} Following the multi-claim convention for this source, source\_url is the source-level URL https://stacks.math.columbia.edu/ and the source label lemma-relatively-very-ample is carried in claim\_key and here rather than in a per-claim tag URL.\par\medskip
\noindent\textit{Editorial note.} The preamble is the chapter's shared Stacks Project preamble, retained as sources/stacks-morphisms/retained-preamble.tex: a line-for-line excerpt of the shared Stacks preamble staged with this snapshot. No macro was invented or redefined, and this card adds no preamble text of its own.\par\medskip
\noindent\textit{Editorial note.} Statement span 9145-9169 and proof span 9171-9237 come from the halfA pool card operations/pool-cards/stacks-morphisms-halfA.json and were re-read in the pinned chapter file before this card was written.\par\medskip
\noindent\textit{Editorial note.} The author writes 'we urge the reader to prove this for themselves' at the point where immersiveness has to be shown, and later 'We omit the verification.' after describing the factorization of the associated morphism through the open subscheme. Both sentences are kept verbatim; the counted span nevertheless contains the reduction, the construction of the open subscheme U\_j and the factorization, so what is left to the reader is the verification that the displayed factorization exhibits an immersion, not the mechanism of the argument.\par\medskip
\noindent\textit{Editorial note.} 5 cross-references occur inside the counted spans; they are listed in external\_references as navigation pointers only.\par\medskip
\noindent\textit{Allowed standard prerequisites (editorial).}\begin{itemize}
\item A relatively very ample sheaf makes the morphism separated, cited from the same chapter, together with locality of separatedness on the base. Cited at proof lines 9181-9183. Gives separatedness needed for the pushforward of the sheaf to be quasi-coherent and for the formulation of the fourth condition.
\item Pushforward along a quasi-compact separated morphism preserves quasi-coherence, cited from the schemes chapter. Cited at proof lines 9169-9171. Makes the pushforward of the invertible sheaf a quasi-coherent module so that its projective bundle and the associated morphism make sense.
\item The description of a morphism into a projective bundle by a surjection from a pullback of a module, and the adjunction with a map into the pushforward, cited from the constructions chapter. Cited at proof lines 9186-9190. Converts the local immersions into projective bundles into the surjectivity statement of the fourth condition.
\item The open subscheme of morphisms of relative Proj's induced by a map of modules, and the closed immersion of relative Proj's induced by a surjection from a symmetric algebra generated in degree one. Cited at proof lines 9200-9203 and 9229-9233. Supplies the open U\_j through which the associated morphism factors, and the final implication from the third condition to the first.
\end{itemize}\medskip
\ifdefined\hypertarget\hypertarget{sourcepointer-1}{}\fi
\noindent\textit{Source reference.} Original lemma 07ZG: focused-seeds/AGP/sources/stacks-snapshot-20261011/tags.txt, source lines 10146--10146. Same-chapter reference used in the retained author text at source line(s) 9208. It is supplied only as a navigation pointer so that the retained text compiles; the counted item delivers the author's own statement and the author's complete proof, and no part of the referenced item is copied in or substituted for the counted argument.\par\medskip
\ifdefined\hypertarget\hypertarget{sourcepointer-2}{}\fi
\noindent\textit{Source reference.} Original lemma 07ZK: focused-seeds/AGP/sources/stacks-snapshot-20261011/tags.txt, source lines 10150--10150. Same-chapter reference used in the retained author text at source line(s) 9074, 9235, 10528. It is supplied only as a navigation pointer so that the retained text compiles; the counted item delivers the author's own statement and the author's complete proof, and no part of the referenced item is copied in or substituted for the counted argument.\par\medskip
\ifdefined\hypertarget\hypertarget{sourcepointer-3}{}\fi
\noindent\textit{Source reference.} Original section 01OA: focused-seeds/AGP/sources/stacks-snapshot-20261011/tags.txt, source lines 2187--2187. Same-chapter reference used in the retained author text at source line(s) 9189. It is supplied only as a navigation pointer so that the retained text compiles; the counted item delivers the author's own statement and the author's complete proof, and no part of the referenced item is copied in or substituted for the counted argument.\par\medskip
\ifdefined\hypertarget\hypertarget{sourcepointer-4}{}\fi
\noindent\textit{Source reference.} Original lemma 01VQ: focused-seeds/AGP/sources/stacks-snapshot-20261011/tags.txt, source lines 2455--2455. Same-chapter reference used in the retained author text at source line(s) 9182. It is supplied only as a navigation pointer so that the retained text compiles; the counted item delivers the author's own statement and the author's complete proof, and no part of the referenced item is copied in or substituted for the counted argument.\par\medskip
\ifdefined\hypertarget\hypertarget{sourcepointer-5}{}\fi
\noindent\textit{Source reference.} Original lemma 01LC: focused-seeds/AGP/sources/stacks-snapshot-20261011/tags.txt, source lines 2081--2081. Same-chapter reference used in the retained author text at source line(s) 262, 404, 882, 1765, 1924, 2256, 2296, 8726, 9175, 10731, 13903. It is supplied only as a navigation pointer so that the retained text compiles; the counted item delivers the author's own statement and the author's complete proof, and no part of the referenced item is copied in or substituted for the counted argument.\par\medskip
\begin{lemma}
\label{lemma-relatively-very-ample}
Let $f : X \to S$ be a morphism of schemes.
Let $\mathcal{L}$ be an invertible sheaf on $X$.
Assume $f$ is quasi-compact. The following are
equivalent
\begin{enumerate}
\item $\mathcal{L}$ is relatively very ample on $X/S$,
\item there exists an open covering $S = \bigcup V_j$ such
that $\mathcal{L}|_{f^{-1}(V_j)}$ is relatively very ample
on $f^{-1}(V_j)/V_j$ for all $j$,
\item there exists a quasi-coherent sheaf of graded
$\mathcal{O}_S$-algebras $\mathcal{A}$ generated in degree
$1$ over $\mathcal{O}_S$ and a map of graded $\mathcal{O}_X$-algebras
$\psi : f^*\mathcal{A} \to \bigoplus_{n \geq 0} \mathcal{L}^{\otimes n}$
such that $f^*\mathcal{A}_1 \to \mathcal{L}$ is surjective and the
associated morphism
$r_{\mathcal{L}, \psi} : X \to \underline{\text{Proj}}_S(\mathcal{A})$
is an immersion, and
\item $f$ is quasi-separated, the canonical map
$\psi : f^*f_*\mathcal{L} \to \mathcal{L}$ is surjective, and
the associated map $r_{\mathcal{L}, \psi} : X \to \mathbf{P}(f_*\mathcal{L})$
is an immersion.
\end{enumerate}
\end{lemma}

\begin{proof}
It is clear that (1) implies (2). It is also clear that
(4) implies (1); the hypothesis of quasi-separation
in (4) is used to guarantee that $f_*\mathcal{L}$ is quasi-coherent via
Schemes, Lemma \ref{schemes-lemma-push-forward-quasi-coherent}.

\medskip\noindent
Assume (2). We will prove (4).
Let $S = \bigcup V_j$ be an open covering as in (2).
Set $X_j = f^{-1}(V_j)$ and $f_j : X_j \to V_j$ the
restriction of $f$. We see that $f$ is separated by
Lemma \ref{lemma-relatively-very-ample-separated} (as being
separated is local on the base). By assumption there exists a
quasi-coherent $\mathcal{O}_{V_j}$-module $\mathcal{E}_j$ and an immersion
$i_j : X_j \to \mathbf{P}(\mathcal{E}_j)$ with
$\mathcal{L}|_{X_j} \cong i_j^*\mathcal{O}_{\mathbf{P}(\mathcal{E}_j)}(1)$.
The morphism $i_j$ corresponds to a surjection
$f_j^*\mathcal{E}_j \to \mathcal{L}|_{X_j}$, see
Constructions, Section \ref{constructions-section-projective-bundle}.
This map is adjoint to a map $\mathcal{E}_j \to f_*\mathcal{L}|_{V_j}$
such that the composition
$$
f_j^*\mathcal{E}_j \to (f^*f_*\mathcal{L})|_{X_j} \to \mathcal{L}|_{X_j}
$$
is surjective. We conclude that
$\psi : f^*f_*\mathcal{L} \to \mathcal{L}$ is surjective. Let
$r_{\mathcal{L}, \psi} : X \to \mathbf{P}(f_*\mathcal{L})$ be the
associated morphism. We still have to show that $r_{\mathcal{L}, \psi}$
is an immersion; we urge the reader to prove this for themselves.
The $\mathcal{O}_{V_j}$-module map $\mathcal{E}_j \to f_*\mathcal{L}|_{V_j}$
determines a homomorphism on symmetric algebras, which in turn defines a
morphism
$$
\mathbf{P}(f_*\mathcal{L}|_{V_j}) \supset U_j \longrightarrow
\mathbf{P}(\mathcal{E}_j)
$$
where $U_j$ is the open subscheme of
Constructions, Lemma \ref{constructions-lemma-morphism-relative-proj}.
The compatibility of $\psi$ with $\mathcal{E}_j \to f_*\mathcal{L}|_{V_j}$
shows that $r_{\mathcal{L}, \psi}(X_j) \subset U_j$ and that there is a
factorization
$$
\xymatrix{
X_j \ar[r]^-{r_{\mathcal{L}, \psi}} & U_j \ar[r] &
\mathbf{P}(\mathcal{E}_j)
}
$$
We omit the verification.
This shows that $r_{\mathcal{L}, \psi}$ is an immersion.

\medskip\noindent
At this point we see that (1), (2) and (4) are equivalent.
Clearly (4) implies (3). Assume (3). We will prove (1).
Let $\mathcal{A}$ be a quasi-coherent sheaf of graded $\mathcal{O}_S$-algebras
generated in degree $1$ over $\mathcal{O}_S$. Consider the map of
graded $\mathcal{O}_S$-algebras $\text{Sym}(\mathcal{A}_1) \to \mathcal{A}$.
This is surjective by hypothesis and hence induces a closed immersion
$$
\underline{\text{Proj}}_S(\mathcal{A})
\longrightarrow
\mathbf{P}(\mathcal{A}_1)
$$
which pulls back $\mathcal{O}(1)$ to $\mathcal{O}(1)$,
see Constructions, Lemma
\ref{constructions-lemma-surjective-generated-degree-1-map-relative-proj}.
Hence it is clear that (3) implies (1).
\end{proof}

\end{document}
